%%%%%%%%%%%%%%%%%%%%%%%%%%%%%
% Best Base Choice Begin here
%%%%%%%%%%%%%%%%%%%%%%%%%%%%%
\subsection{Optimal Base Size Choice}
The base multiplier size strongly affects DSP usage, hierarchy depth, and latency of large multipliers.
We evaluate candidate bases produced in Section~\ref{sec_build_base_mult}.
% Figure~\ref{fig_dsp_waste_vs_mult_size} plots DSP waste above the theoretical minimum from Equation~\ref{eq_mul_best}.
Bases derived by fixing the larger DSP dimension ($26$ for DSP48E2 $26\times17$) show consistently lower waste (Equation~\ref{eq_mul_best}).
Among those candidates, the $77\times77$ base achieves the lowest waste across many sizes.
Low DSP waste alone does not guarantee an efficient hierarchy for Karatsuba or Schoolbook decomposition.
We therefore consider nearby sizes that improve hierarchical coverage, even if they generate slightly larger than the target.
For example, tiling $77\times77$ toward a $1024\times1024$ target yields a largest hierarchical block of $924\times924$ and leaves two non-hierarchical rectangular regions.
In contrast, a $43\times43$ base can reach $1032\times1032$ and thus increase the hierarchical nature for Karatsuba.
We capture these trade-offs using a Pareto frontier analysis across targets from $160$ to $1024$ bits.

% % dsp waste graph 
% \begin{figure}[t]
%     \centering
%     \includegraphics[width=0.8\columnwidth]{figures/base_tiling/dsp_waste_vs_multiplier_side.pdf}
%     \caption{DSP waste versus multiplier size for different base multiplier configurations.}
%     \label{fig_dsp_waste_vs_mult_size}
% \end{figure}



The DSP cost of a fully hierarchical tiling is known once the base and the decomposition schemes are chosen.
Let the base size be $B=nW_1+mW_2$ with $W_1=26$ and $W_2=17$ for DSP48E2.
If the hierarchical side length is $T=2^a\times 3^b\times B$, and 
$\text{DSP}_{\text{base}}$ is the DSP count of the base multiplier.
The total DSP count for the hierarchical region is
\begin{equation}
\text{DSP}_{\text{total}} = \text{DSP}_{\text{base}} \times N_{\text{base}}
\end{equation}

where $N_{\text{base}}$ is the number of base multipliers implied by the hierarchy.

\begin{equation}
N_{\text{base}} = 
\begin{cases}
4^a \times 9^b & \text{(Schoolbook at all levels)} \\
3^a \times 6^b & \text{(Karatsuba at all levels)} \\
\text{Mixed} & \text{(Combination of schemes)}
\end{cases}
\end{equation}


Figure~\ref{fig_pareto_frontier_all_configs} shows the Pareto frontier across decomposition schedules and highlights the best base per target size.
Across the evaluated targets, $43\times43$ minimizes DSP usage when Karatsuba is preferred.
Across the evaluated targets, $128\times128$ is often best when Schoolbook is preferred.




% create two figures indicating pareto font

\begin{figure}[t]
    \centering
    \includegraphics[width=\columnwidth]{figures/base_tiling/pareto_frontier_best_bases.pdf}
    \caption{
        Pareto frontier for all configurations.
        The annotation shows the achieved size using the best base per target.
    }
    \label{fig_pareto_frontier_all_configs}
\end{figure}









%%%%%%%%%%%%%%%%%%%%%%%%%%%%%%
% Tiler Begin here
%%%%%%%%%%%%%%%%%%%%%%%%%%%%%%
\subsection{Tiler} \label{sec_tiler}
Finding the optimal placement of small rectangles to tile a large rectangle is NP-complete~\cite{BeumontTilingNPComplete2002}.
General tiling schemes suffers from two main limitations when applied to FPGA-based large multipliers:
\begin{itemize}
    \item They do not prioritize hierarchical regions that enable Karatsuba decomposition.
    \item They ignore architecture specific features such as DSP cascade capabilities and efficient base sizes.
\end{itemize}


We evaluate two tiling methodologies.
The first method tiles to the exact target size.
We illustrate the exact tiling flow on a $1024\times1024$ multiplier using a $77\times77$ base.
First step, the algorithm selects the largest hierarchical side length of the form $H=2^x\times3^y\times B$ that is less than or equal to the target.
For $T=1024$ and $B=77$, the largest hierarchical region is $H=924$ with $x=2$ and $y=1$.
This yields a $12\times12$ grid of base multipliers with hierarchy of a $3\times3$ grid tiled in a $2\times2$ fashion two times.
This leaves two rectangle areas around the hierarchical region that must be tiled further.
We explore three variants for tiling the remaining area.
The first variant is greedily placing the same base until it does not fit, then placing DSP-fit multiplications.
The second variant recursively tiles the remaining area using other efficient bases (e.g., $44\times44$ and $94\times94$) and small DSP-fit tiles.
The third variant improves on this by maximizing the number of cascade-compatible connections.
The cascade-compatible variant yields the best DSP usage for the remaining area among the three variants($464$ compared to $516$ in the first variant and $507$ in the second variant).
In the three variants, each candidate placement is scored by DSP count and by estimated LUT/FF overhead for accumulation.
Cascade-compatible connections are preferred because they reduce LUT-based shifting and addition operations.
Figure~\ref{fig_tiling1024_recurmax} shows the recursive approach without cascade consideration, the rest is omitted for brevity.
Then we compare to the second method (non-exact tiling) that may tile to slightly larger than the target to maximize the generated hierarchical area, then we choose the best method of the two.


\begin{figure}[t]
     \centering
     \includegraphics[width=\columnwidth]{figures/tiling/tile1024RecurMax.pdf}
     \caption{Recursive tiling approach without cascade consideration for a $1024\times1024$ multiplier using $516$ DSP48E2 slices. The purple boxes highlight remaining rectangular  regions after hierarchical tiling.}
     \label{fig_tiling1024_recurmax}
\end{figure}


% \begin{figure}[t]
%      \centering
%     %  \begin{subfigure}{0.475\columnwidth}
%      \begin{subfigure}{0.9\columnwidth}
%          \centering
%          \includegraphics[width=\textwidth]{figures/tiling/tile1024Greedy.pdf}
%          \caption{Greedy approach.}
%          \label{fig_tile1024_greedy}
%      \end{subfigure}
%      \hfill

%     %  \begin{subfigure}{0.475\columnwidth}
%      \begin{subfigure}{0.9\columnwidth}
%          \centering
%          \includegraphics[width=\textwidth]{figures/tiling/tile1024RecurMax.pdf}
%          \caption{Recursive approach without cascade consideration.}
%          \label{fig_tiling1024_recurmax}
%      \end{subfigure}
%      \hfill 

%     %  \begin{subfigure}{0.475\columnwidth}
%      \begin{subfigure}{0.9\columnwidth}
%          \centering
%          \includegraphics[width=\textwidth]{figures/tiling/tile1024max17.pdf}
%          \caption{Recursive approach with cascade consideration.}
%          \label{fig_tiling1024_max}
%      \end{subfigure}
%      \hfill 


%     \caption{
%         Different tiling schemes for a $1024\times1024$ multiplier.
%         The greedy approach uses $507$ DSP48E2 slices.
%         The recursive approach without cascade uses $516$ DSP48E2 slices.
%         The recursive approach with cascade uses $464$ DSP48E2 slices.
%     }
%     \label{fig_tiling1024x1024}
    
% \end{figure}

%%%%%%%%%%%%%%%%%%%%%%%%%%%%%%
% Tiler end here
%%%%%%%%%%%%%%%%%%%%%%%%%%%%%%